\chapter{Related Work}
The review covers relatively latest works in the areas of vision--language model design,
text-rich and document-heavy images, charts and structured visual reasoning, faithfulness and critique, and efficient adaptation.
We group it by the parts that actually matter here the most: model design, text-rich images, chart reasoning, faithfulness, datasets, metrics, and low-budget adaptation.

\section{Recent vision--language model design}
Recent vision--language models are not built around a single alignment trick.
Their performance is a result of many choices working together, such as the backbone architecture,
the visual tokenisation method, the data mixture, the connector design, and the instruction tuning.
Since all those factors are usually changed together, it is hard to tell which part actually caused an improvement in the final performance.

This is exactly the issue raised by recent ablation studies, 
which demonstrate that these factors are still poorly separated, and that it is vital to conduct controlled comparisons~\cite{laurencon2024whatmatters}.
MM1.5 also makes the same argument, showing that optical character recognition (OCR) data,
synthetic captions, instruction mixtures, and specialised variants 
can each weigh as much as raw scale in the final performance of a vision--language model~\cite{zhang2025mm15}.
Molmo and PixMo are useful here because they provide open weights and open data 
that let us actually read the recipe and understand what is going on under the hood~\cite{deitke2025molmo}.
LLaVA-OneVision shows that a single family of models can be used across image and video when the data mixture is chosen with care~\cite{li2025llavaonevision}.
The same development can also be seen in earlier and related open VLM families.
OFA and GIT frame visual tasks as generation problems,
while Qwen-VL, PaLI, OpenFlamingo, CogVLM, LLaVA, MiniGPT-4, and InstructBLIP show different ways of connecting a visual encoder with a language model and instruction-tuning it for practical use~\cite{wang2022ofa,wang2022git,bai2023qwenvl,chen2023pali,awadalla2023openflamingo,wang2024cogvlm,liu2023llava,zhu2023minigpt4,dai2023instructblip}.

For this study, the conclusion is simple.
If the model, the data, and the budget are all changing together, then the contribution of the objective itself is hard to tell.
That is why we need a controlled comparison, where the model, the data, and the budget are held constant while only the training signal changes.

Open model families are also important, because they enable us to actually run the controlled comparison on a practical budget.
Qwen2-VL is designed to work with dynamic image resolution and text-heavy visual inputs such as documents, charts, and diagrams~\cite{wang2024qwen2vl}.
BLIP-2's Q-Former lets us test an image--answer contrastive loss while keeping the larger vision encoder and language model frozen~\cite{li2023blip2}.
We use these open models because they make controlled, parameter-efficient experiments feasible under limited resources.

\section{Text-rich and document-heavy images}
Documents are difficult for a specific reason.
The answer in this case is not in the whole image, but rather in a small part of it, such as a line of text, a table cell, or a layout region.
Hence, the task is not just an image recognition problem.
The model has to read the document and find the evidence in the right place, which is a different skill than naming objects in a natural image.

Recent work shows that text-rich image reasoning is moving toward multimodal large language models, especially for OCR-heavy and layout-heavy inputs~\cite{fu2025textrichsurvey}.
DocLLM is a good example of this trend. It encodes the document layout directly into the generative document-understanding problem,
which is a step toward more faithful reasoning on documents~\cite{wang2024docllm}.
OmniDocBench focuses on the coverage angle instead, scoring diverse PDF parsing cases with detailed annotations,
which is a step toward more structured evaluation of document understanding~\cite{ouyang2025omnidocbench}.
In this study, we use DocVQA as a test case for document understanding, because it is a well-known benchmark 
where the answer is usually located on scanned text, layout, and local visual context~\cite{mathew2021docvqa}.

Our goal is not to build a new document model.
We want to see how much the training objective itself contributes to the model's ability to read documents and answer questions about them.

\section{Charts and structured visual reasoning}
Charts create a different kind of problem.
In this case, the model has to read the labels, understand the axes and tick marks,
and determine the relationships between the visual elements, such as bars, lines, or points.
ChartInstruct and ChartGemma demonstrate that instruction tuning on chart images can help in both chart understanding and reasoning~\cite{masry2024chartinstruct,masry2025chartgemma}.
ChartMimic shows the same point by training a model to generate code for that chart, which makes the model understand the chart in a more structured way~\cite{yang2025chartmimic}.
ChartMuseum shows that current large visual--language models still perform far poorer than humans on chart reasoning~\cite{tang2025chartmuseum}.
MathVista, MMMU-Pro, and MMReason extend this concern to mathematical and multi-step reasoning, which is often required for chart understanding~\cite{lu2024mathvista,yue2025mmmupro,yao2025mmreason}.

This is why we use ChartQA, ScienceQA, and A-OKVQA in the same study~\cite{masry2022chartqa,lu2022scienceqa,schwenk2022aokvqa}.
Together, they cover structured chart reasoning, science reasoning, and knowledge-heavy visual reasoning.


\section{Faithfulness, critique, and efficient adaptation}
Explanation-aware training has a subtle problem.
The model might learn to write an explanation that looks convincing,
but that is not actually based on the image and is not the real cause leading to the answer.
Recent studies approach this problem from different angles.
MMMU tests expert-level multimodal reasoning across many visual formats, to see if VLMs can do expert-style visual reasoning and not just simple image captioning~\cite{yue2024mmmu}.
MMStar checks whether a model actually needs the image to answer the question or whether the answer can be derived from the learned knowledge alone~\cite{chen2024mmstar}.
HallusionBench focuses on cases where language hallucination and visual illusion appear together in LVLM responses~\cite{guan2024hallusionbench}.
Hallucination-augmented contrastive learning turns a hallucination into a contrastive signal, to teach the model to avoid it~\cite{jiang2024hacl}.
VISCO and Critic-V evaluate the model's ability to critique, correct, and detect errors in its own output~\cite{wu2025visco,zhang2025criticv}.
General multimodal benchmarks such as POPE, MME, MMBench, SEED-Bench, and MM-Vet are also relevant here,
because they test hallucination, perception, reasoning, and instruction-following outside one narrow dataset score~\cite{li2023pope,fu2023mme,liu2023mmbench,li2024seedbench,yu2024mmvet}.
These works suggest that faithfulness should be measured separately from answer accuracy.

Efficient adaptation is also important,
because it allows us to test the effect of the training objective without needing to fine-tune the entire model.
DoRA is a good example of this approach, which demonstrates that we can actually adapt the model by updating only a small portion of the parameters~\cite{liu2024dora}.
Reason-RFT points to another direction that is relevant here, where post-training specifically rewards visual reasoning~\cite{tan2025reasonrft}.
In this paper, we stay with supervised objective comparisons, but the concern is similar.
We evaluate the model as a reasoning system, rather than only as an answer generator.

Table~\ref{tab:LiteratureSummary} summarises the relevant works and the gaps we address in this study.
It is not meant to be a comprehensive review of the field, but rather a snapshot of the recent trends and how they relate to our controlled comparison.

\begin{table*}[!ht]
\centering
\caption{Representative related work and how it connects to the gap studied here. A check mark means the line of work directly covers the property, a cross means it is not part of the main scope, and a tilde means partial coverage.}
\label{tab:LiteratureSummary}
\scriptsize
\begin{adjustbox}{max width=\textwidth}
\begin{tabular}{|p{3.2cm}|p{2.4cm}|c|c|c|c|p{4.5cm}|}
\hline
\textbf{Work} & \textbf{Main focus} & \textbf{Generation} & \textbf{Structured VQA} & \textbf{Rationales} & \textbf{Efficient FT} & \textbf{Connection to this study} \\ \hline
Recent VLM construction~\cite{laurencon2024whatmatters,deitke2025molmo,li2025llavaonevision,zhang2025mm15} & Model and data recipe & \cmark & \pmark & \pmark & \pmark & Shows that model design and training data must be controlled before we can attribute gains to the objective. \\ \hline
Open VLM backbones~\cite{wang2024qwen2vl,li2023blip2} & Adaptable model families & \cmark & \pmark & \pmark & \pmark & Provide practical backbones for controlled experiments under a fixed compute budget. \\ \hline
Text-rich document work~\cite{fu2025textrichsurvey,wang2024docllm,ouyang2025omnidocbench} & Documents and PDF parsing & \cmark & \cmark & \pmark & \pmark & Motivates the document setting, but usually does not compare training objectives under one fixed backbone. \\ \hline
Chart reasoning work~\cite{masry2024chartinstruct,masry2025chartgemma,tang2025chartmuseum,yang2025chartmimic} & Chart understanding & \cmark & \cmark & \pmark & \pmark & Shows that chart reasoning is still difficult and needs structured evaluation beyond ordinary image QA. \\ \hline
Recent reasoning benchmarks~\cite{lu2024mathvista,yue2024mmmu,chen2024mmstar,yue2025mmmupro,yao2025mmreason} & Harder multimodal evaluation & \pmark & \cmark & \pmark & \xmark & Motivate evaluation on visual, mathematical, expert-level, and multi-step reasoning tasks. \\ \hline
Faithfulness and critique~\cite{guan2024hallusionbench,jiang2024hacl,wu2025visco,zhang2025criticv} & Reliability of reasoning & \pmark & \cmark & \cmark & \xmark & Support the need to measure hallucination, critique, and evidence use separately from answer accuracy. \\ \hline
DoRA / Reason-RFT~\cite{liu2024dora,tan2025reasonrft} & Efficient and reasoning-focused adaptation & \pmark & \xmark & \pmark & \cmark & Motivate low-resource adaptation and explicit reasoning-oriented post-training. \\ \hline
This study & Controlled explanation-aware alignment & \cmark & \cmark & \cmark & \cmark & Compares objective choices directly and reports answer accuracy, rationale quality, and evidence sensitivity together. \\ \hline
\end{tabular}
\end{adjustbox}
\end{table*}

A common pattern emerges from this review.
Earlier work on multimodal alignment focused mainly on learning image--text representations, or on generating an answer from the image and the question.
More recent work has started to pay attention to the faithfulness of the model's reasoning, and to the quality of the explanations it generates.
However, it is still unclear how much the training objective itself contributes to the model's performance.
In this study, our goal is to fill that gap by conducting a controlled comparison of different training objectives while keeping all other factors constant.
